\begin{abstract}
A model that reads an updated record can report the current value correctly and still let the old value
affect what it decides. We test whether that happens, and whether it happens \emph{because} the old value
was an assignment that has since been superseded. In synthetic access-control and routing records, the
current state, the decision policy and the correct action are held fixed while a single obsolete state is
edited so that it implies either the correct or the wrong action. Matched controls place the same word as
an assignment to an unrelated attribute that was later updated, as a word the same entity merely mentioned,
or as an unattached word. In a frozen confirmation on held-out wording (Qwen3-8B, 768 histories), the
obsolete state raised the wrong action's log-odds by 3.904 and 3.133 nats in the two tasks and flipped the
generated decision in about one history in six, almost always toward the wrong action (65 against 1 and 75
against 4 histories), including in routing conditions where the model retrieved both the current
(0.982--0.988) and the initial (0.990) state and decided correctly without history (0.953--0.961). The
effect is not specific to supersession: an incidental mention of the same word by the same entity moved
decisions as much in access and more in routing (net switches 0.294 against 0.185), and whether a
superseded assignment exceeded the mention control in score depended on the prompt template. A superseded
assignment exceeded an updated unrelated attribute in routing but not detectably in access. Finally,
the edit's effect on easy retrieval queries, a score-level sensitivity measure, added no information
about which decisions would fail beyond simple pre-decision predictors (log-loss gain $+0.0004$ nats,
95\% interval $[-0.0011, +0.0019]$). A second model, Gemma 3 4B, did not meet the task-competence gate,
so the results are for one model.
\end{abstract}
